% NGÂN HÀNG TUYỂN SINH 10 – BÀI 4

\begin{bt}%[TS10 4.8]%[ma:0a19aad1]
% nguon: Đề tham khảo Tuyển sinh 10 TP.HCM 2025-2026 – Quận 3 – Đề 1 – Bài 4
Một tổ may gồm cả công nhân nam và nữ được giao may $360$ chiếc áo cổ vũ. Để hoàn thành nhiệm vụ, mỗi công nhân nam may $6$ chiếc áo, mỗi công nhân nữ may $8$ chiếc áo. Gọi $x$ và $y$ lần lượt là số công nhân nam và số công nhân nữ ($x, y \in \mathbb{N}^*$).
\begin{enumerate}[a)]
	\item Lập công thức biểu diễn $y$ theo $x$.
	\item Tổ may có bao nhiêu công nhân? Biết rằng số công nhân nữ gấp $1{,}5$ lần số công nhân nam.
\end{enumerate}
\loigiai{
\begin{enumerate}[a)]
	\item Vì $x$ công nhân nam may được $6x$ chiếc áo và $y$ công nhân nữ may được $8y$ chiếc áo, tổng số áo may được là $360$ chiếc nên ta có phương trình:
	$6x + 8y = 360 \Leftrightarrow 8y = 360 - 6x \Leftrightarrow y = \dfrac{360 - 6x}{8} = 45 - \dfrac{3}{4}x$.
	\item Theo giả thiết, số công nhân nữ gấp $1{,}5$ lần số công nhân nam, tức là:
	$y = 1{,}5x = \dfrac{3}{2}x$.\\
	Thay vào phương trình $6x + 8y = 360$, ta được:
	$6x + 8\left(\dfrac{3}{2}x\right) = 360 \Leftrightarrow 6x + 12x = 360 \Leftrightarrow 18x = 360 \Leftrightarrow x = 20\text{ (công nhân)}$.\\
	Suy ra số công nhân nữ là: $y = 1{,}5 \cdot 20 = 30\text{ (công nhân)}$.\\
	Vậy tổ may có tất cả $20 + 30 = 50\text{ công nhân}$.
\end{enumerate}
}
\end{bt}

\begin{bt}%[TS10 4.5]%[ma:4ef0dfb7]
% nguon: Đề tham khảo Tuyển sinh 10 TP.HCM 2025-2026 – Quận 4 – Đề 1 – Bài 4
Một câu lạc bộ bóng rổ tổ chức tuyển chọn thành viên mới với quy định: Mỗi học sinh thực hiện $25$ lượt ném bóng. Mỗi lượt ném bóng vào rổ được cộng $5$ điểm, mỗi lượt ném không vào rổ bị trừ $2$ điểm. Ban đầu mỗi bạn có sẵn $50$ điểm. Để được chọn vào đội thì sau $25$ lượt ném, học sinh phải đạt tổng số điểm từ $100$ điểm trở lên.
\begin{enumerate}[a)]
	\item Một bạn học sinh thực hiện $25$ lượt ném gồm $18$ lượt bóng vào rổ và $7$ lượt bóng không vào rổ. Tính tổng số điểm của học sinh đó.
	\item Gọi $x$ là số lượt ném bóng vào rổ trong $25$ lượt ($x \in \mathbb{N}, 0 \le x \le 25$). Hãy biểu diễn tổng số điểm của học sinh theo $x$. Hỏi để được chọn vào đội, học sinh cần ném bóng vào rổ ít nhất bao nhiêu lượt?
\end{enumerate}
\loigiai{
\begin{enumerate}[a)]
	\item Tổng số điểm của học sinh sau khi thực hiện $18$ lượt vào rổ và $7$ lượt không vào rổ là:
	$50 + 18 \cdot 5 - 7 \cdot 2 = 50 + 90 - 14 = 126\text{ (điểm)}$.
	\item Trong $25$ lượt ném, nếu có $x$ lượt bóng vào rổ thì số lượt bóng không vào rổ là $25 - x$.\\
	Biểu thức biểu diễn tổng số điểm của học sinh theo $x$ là:
	$T(x) = 50 + 5x - 2(25 - x) = 50 + 5x - 50 + 2x = 7x\text{ (điểm)}$.\\
	Để học sinh được chọn vào đội, tổng số điểm phải đạt từ $100$ điểm trở lên:
	$7x \ge 100 \Leftrightarrow x \ge \dfrac{100}{7} \approx 14{,}28$.\\
	Vì $x$ là số nguyên dương và $x \le 25$ nên $x \ge 15$.\\
	Vậy học sinh cần thực hiện ít nhất $15$ lượt ném bóng vào rổ thì mới được chọn vào đội.
\end{enumerate}
}
\end{bt}

\begin{bt}%[TS10 4.3]%[ma:a3ef1068]
% nguon: Đề tham khảo Tuyển sinh 10 TP.HCM 2025-2026 – Quận 4 – Đề 2 – Bài 4
Bạn Nam mang theo $200\,000\text{ đồng}$ vào nhà sách mua đồ dùng học tập. Nhà sách đang có chương trình khuyến mãi giảm giá $10\%$ so với giá niêm yết cho toàn bộ các mặt hàng. Biết giá niêm yết của bút viết là $5000\text{ đồng/cây}$, tập vở là $10\,000\text{ đồng/cuốn}$, thước kẻ là $6000\text{ đồng/cây}$. Bạn Nam mua $x$ cây viết, $6$ quyển tập và $4$ cây thước kẻ. Gọi $y$ (đồng) là số tiền còn lại của bạn Nam sau khi thanh toán.
\begin{enumerate}[a)]
	\item Hãy biểu diễn $y$ theo $x$ (rút gọn biểu thức).
	\item Hỏi bạn Nam có thể mua được nhiều nhất bao nhiêu cây viết?
\end{enumerate}
\loigiai{
\begin{enumerate}[a)]
	\item Giá bán sau khi giảm giá $10\%$ của từng loại mặt hàng là:
	\begin{itemize}
		\item Giá một cây viết: $5000 \cdot (1 - 0{,}1) = 4500\text{ (đồng/cây)}$.
		\item Giá một quyển tập: $10\,000 \cdot (1 - 0{,}1) = 9000\text{ (đồng/cuốn)}$.
		\item Giá một cây thước: $6000 \cdot (1 - 0{,}1) = 5400\text{ (đồng/cây)}$.
	\end{itemize}
	Tổng số tiền Nam mua $6$ quyển tập và $4$ cây thước là:
	$6 \cdot 9000 + 4 \cdot 5400 = 54\,000 + 21\,600 = 75\,600\text{ (đồng)}$.\\
	Tổng số tiền Nam mua $x$ cây viết là $4500x\text{ (đồng)}$.\\
	Biểu thức số tiền còn lại của bạn Nam là:
	$y = 200\,000 - (75\,600 + 4500x) = 124\,400 - 4500x\text{ (đồng)}$.
	\item Để Nam có đủ tiền thanh toán thì số tiền còn lại $y \ge 0$:\\
	$124\,400 - 4500x \ge 0 \Leftrightarrow 4500x \le 124\,400 \Leftrightarrow x \le \dfrac{124\,400}{4500} \approx 27{,}64$.\\
	Vì số cây viết $x$ là số nguyên không âm nên $x \le 27$.\\
	Vậy bạn Nam có thể mua được nhiều nhất $27$ cây viết.
\end{enumerate}
}
\end{bt}

\begin{bt}%[TS10 4.5]%[ma:4dc125a1]
% nguon: Đề tham khảo Tuyển sinh 10 TP.HCM 2025-2026 – Quận 5 – Đề 2 – Bài 4
Một xí nghiệp dự tính vận chuyển hàng bằng hai xe tải và đang phân vân giữa hai phương án là mua hẳn hai chiếc xe tải hoặc thuê hai xe tải (hai xe giống nhau). Nếu mua hai xe với giá mỗi xe là $200\,000\,000$ đồng thì mỗi ngày xí nghiệp phải trả tiền công cho tài xế và mua nhiên liệu cho một xe là $5\,000\,000$ đồng. Còn nếu thuê xe thì giá thuê một xe tải là $10\,000\,000$ đồng/ngày (đã bao gồm tiền công cho tài xế và nhiên liệu).
\begin{enumerate}[a)]
    \item Gọi $T$ là tổng số tiền xí nghiệp bỏ ra để vận chuyển hàng sau $n$ ngày. Lập công thức tính $T$ theo $n$ đối với mỗi phương án.
    \item Hỏi xí nghiệp vận chuyển hàng từ bao nhiêu ngày trở lên thì nên chọn phương án mua xe sẽ tiết kiệm hơn phương án thuê xe?
\end{enumerate}
\loigiai{
\begin{enumerate}[a)]
    \item Lập công thức tính tổng chi phí $T$ sau $n$ ngày vận chuyển:
    \begin{itemize}
        \item Phương án mua 2 xe tải:\\
        Chi phí mua 2 xe tải ban đầu là:
        $2 \cdot 200\,000\,000 = 400\,000\,000$ (đồng).\\
        Chi phí tài xế và nhiên liệu cho cả 2 xe mỗi ngày là:
        $2 \cdot 5\,000\,000 = 10\,000\,000$ (đồng/ngày).\\
        Tổng số tiền chi sau $n$ ngày theo phương án mua xe là:
        $T_1(n) = 400\,000\,000 + 10\,000\,000n$ (đồng).
        \item Phương án thuê 2 xe tải:\\
        Giá thuê cả 2 xe mỗi ngày là:
        $2 \cdot 10\,000\,000 = 20\,000\,000$ (đồng/ngày).\\
        Tổng số tiền chi sau $n$ ngày theo phương án thuê xe là:
        $T_2(n) = 20\,000\,000n$ (đồng).
    \end{itemize}
    \item Để phương án mua xe tiết kiệm hơn phương án thuê xe thì tổng chi phí mua xe phải nhỏ hơn tổng chi phí thuê xe:
    $T_1(n) < T_2(n) \Leftrightarrow 400\,000\,000 + 10\,000\,000n < 20\,000\,000n \Leftrightarrow 10\,000\,000n > 400\,000\,000 \Leftrightarrow n > 40$.\\
    Vì $n$ là số ngày nguyên dương nên $n \ge 41$.\\
    Vậy xí nghiệp vận chuyển hàng từ 41 ngày trở lên thì nên chọn phương án mua xe sẽ tiết kiệm hơn phương án thuê xe.
\end{enumerate}
}
\end{bt}

\begin{bt}%[TS10 4.3]%[ma:a9472a44]
% nguon: Đề tham khảo Tuyển sinh 10 năm học 2025-2026 Sở GD&ĐT TP.HCM – TP. Thủ Đức - Đề 3 (Bài 4)
Một siêu thị thực hiện chương trình khuyến mãi nhân dịp lễ: giảm giá $1\,000$ đồng cho mỗi hộp sữa chua nha đam và giảm giá $15\%$ cho mỗi gói kẹo bốn mùa. Một lốc sữa chua nha đam gồm $4$ hộp có giá niêm yết là $30\,000$ đồng/lốc.
\begin{enumerate}[a)]
    \item Bạn Lan mua $5$ lốc sữa chua nha đam thì phải thanh toán tổng cộng bao nhiêu tiền (biết thuế giá trị gia tăng GTGT là $10\%$)?
    \item Bạn An mua $5$ lốc sữa chua nha đam và mua thêm $3$ gói kẹo bốn mùa thì phải thanh toán tất cả $311\,300$ đồng (đã bao gồm thuế GTGT $10\%$). Hỏi giá niêm yết ban đầu của một gói kẹo bốn mùa khi chưa giảm giá là bao nhiêu?
\end{enumerate}
\loigiai{
\begin{enumerate}[a)]
    \item Mỗi lốc sữa chua gồm $4$ hộp nên được giảm:
    $4 \cdot 1\,000 = 4\,000$ (đồng/lốc).
    Giá bán khuyến mãi của một lốc sữa chua khi chưa tính thuế là:
    $30\,000 - 4\,000 = 26\,000$ (đồng).
    Số tiền mua $5$ lốc sữa chua trước thuế là:
    $5 \cdot 26\,000 = 130\,000$ (đồng).
    Số tiền bạn Lan phải trả sau khi tính thuế GTGT $10\%$ là:
    $130\,000 \cdot (1 + 0{,}10) = 130\,000 \cdot 1{,}1 = 143\,000$ (đồng).
    \item Tổng số tiền bạn An đã trả trước khi tính thuế GTGT là:
    $311\,300 : 1{,}1 = 283\,000$ (đồng).
    Số tiền mua $3$ gói kẹo bốn mùa sau khi đã giảm giá là:
    $283\,000 - 130\,000 = 153\,000$ (đồng).
    Giá bán thực tế của một gói kẹo bốn mùa sau khi giảm giá $15\%$ là:
    $153\,000 : 3 = 51\,000$ (đồng).
    Giá niêm yết ban đầu của một gói kẹo bốn mùa là:
    $51\,000 : (1 - 0{,}15) = 51\,000 : 0{,}85 = 60\,000$ (đồng).
\end{enumerate}
}
\end{bt}

\begin{bt}%[TS10 4.5]%[ma:1698752c]
% nguon: Đề tham khảo Tuyển sinh 10 năm học 2025-2026 Sở GD&ĐT TP.HCM – Quận Tân Phú - Đề 1 (Bài 4)
Một khu vui chơi đưa ra hai gói giá vé ưu đãi cho các nhóm gia đình như sau:
\begin{itemize}
    \item Chương trình 1: Giá vé người lớn là $100\,000$ đồng/vé; mỗi vé trẻ em có giá gốc $50\,000$ đồng được giảm $20\,000$ đồng (còn $30\,000$ đồng/vé).
    \item Chương trình 2: Giá vé người lớn là $80\,000$ đồng/vé; giá vé trẻ em giữ nguyên là $50\,000$ đồng/vé.
\end{itemize}
\begin{enumerate}[a)]
    \item Gọi $x$ là số người lớn trong nhóm ($x \in \mathbb{N}^*$). Hãy lập biểu thức biểu thị tổng số tiền vé nhóm phải trả theo Chương trình 1, biết rằng nhóm có $4$ trẻ em.
    \item Hỏi nhóm đó cần có tối thiểu bao nhiêu người lớn để việc lựa chọn mua vé theo Chương trình 2 tiết kiệm chi phí hơn Chương trình 1?
\end{enumerate}
\loigiai{
\begin{enumerate}[a)]
    \item Theo Chương trình 1, giá vé của mỗi trẻ em là:
    $50\,000 - 20\,000 = 30\,000$ (đồng).
    Tổng số tiền vé nhóm gồm $x$ người lớn và $4$ trẻ em phải trả theo Chương trình 1 là:
    $T_1(x) = 100\,000x + 4 \cdot 30\,000 = 100\,000x + 120\,000$ (đồng).
    \item Theo Chương trình 2, tổng số tiền vé nhóm phải trả là:
    $T_2(x) = 80\,000x + 4 \cdot 50\,000 = 80\,000x + 200\,000$ (đồng).
    Để Chương trình 2 có lợi hơn (chi phí thấp hơn) Chương trình 1 thì:
    $T_2(x) < T_1(x) \iff 80\,000x + 200\,000 < 100\,000x + 120\,000 \iff 20\,000x > 80\,000 \iff x > 4$.
    Vì $x$ là số nguyên dương nên giá trị nhỏ nhất của $x$ là $5$.
    Vậy nhóm gia đình cần có tối thiểu $5$ người lớn để chọn Chương trình 2 có lợi hơn.
\end{enumerate}
}
\end{bt}

\begin{bt}%[TS10 4.1]%[ma:686ba9c1]
% nguon: Đề thi chính thức Tuyển sinh 10 TP.HCM năm học 2025-2026 – Bài 4
Một khu vườn hình chữ nhật có chiều dài là $2x\text{ (m)}$ và chiều rộng là $x\text{ (m)}$ ($x > 4$). Bác Ba làm một lối đi xung quanh khu vườn rộng $2\text{ mét}$. Phần đất còn lại ở bên trong dùng để trồng hoa.
\begin{enumerate}[a)]
	\item Viết biểu thức theo $x$ biểu diễn diện tích phần đất dùng để trồng hoa và thu gọn biểu thức đó.
	\item Giả sử diện tích phần đất trồng hoa là $4800\text{ m}^2$. Tính chiều dài và chiều rộng của khu vườn.
\end{enumerate}
\loigiai{
\begin{enumerate}[a)]
	\item Vì lối đi quanh khu vườn rộng $2\text{ m}$ nên phần đất trồng hoa bên trong là hình chữ nhật có:
	\begin{itemize}
		\item Chiều rộng là: $x - 2\cdot 2 = x - 4\text{ (m)}$.
		\item Chiều dài là: $2x - 2\cdot 2 = 2x - 4\text{ (m)}$.
	\end{itemize}
	Biểu thức biểu diễn diện tích phần đất dùng để trồng hoa là:
	$S = (x - 4)(2x - 4) = 2x^2 - 4x - 8x + 16 = 2x^2 - 12x + 16\text{ (m}^2\text{)}$.
	\item Theo giả thiết, diện tích đất trồng hoa là $4800\text{ m}^2$ nên ta có phương trình:
	$2x^2 - 12x + 16 = 4800 \Leftrightarrow 2x^2 - 12x - 4784 = 0 \Leftrightarrow x^2 - 6x - 2392 = 0$.\\
	Ta có $\Delta' = (-3)^2 - 1\cdot(-2392) = 9 + 2392 = 2401 = 49^2 > 0$.\\
	Phương trình có hai nghiệm:
	$x_1 = 3 + 49 = 52$ (thoả mãn điều kiện $x > 4$),\\
	$x_2 = 3 - 49 = -46$ (loại vì $x > 4$).\\
	Với $x = 52$, ta có:
	\begin{itemize}
		\item Chiều rộng của khu vườn là: $x = 52\text{ (m)}$.
		\item Chiều dài của khu vườn là: $2x = 2\cdot 52 = 104\text{ (m)}$.
	\end{itemize}
	Vậy chiều dài của khu vườn là $104\text{ m}$ và chiều rộng là $52\text{ m}$.
\end{enumerate}
}
\end{bt}

\begin{bt}%[TS10 4.1]%[ma:48eaceb3]
% nguon: Đề thi tham khảo Tuyển sinh 10 TP.HCM – Phụ lục Sở GD&ĐT TP.HCM – Bài 4
Một khu vườn hình chữ nhật có chiều dài và chiều rộng lần lượt là $70\text{ m}$ và $30\text{ m}$. Người ta dự tính mở rộng thêm khu vườn bằng cách cải tạo thêm $x\text{ (mét)}$ về mỗi phía ngoài của chiều dài và chiều rộng khu vườn.
\begin{enumerate}[a)]
	\item Viết biểu thức $S$ biểu diễn theo $x$ diện tích của khu vườn hình chữ nhật sau khi mở rộng.
	\item Biết rằng sau khi mở rộng thì diện tích của khu vườn lớn hơn diện tích ban đầu $1150\text{ m}^2$. Tìm giá trị của $x$ (làm tròn đến hàng phần mười của mét).
\end{enumerate}
\loigiai{
\begin{enumerate}[a)]
	\item Sau khi mở rộng thêm $x\text{ (m)}$ về cả hai phía ngoài của chiều dài và chiều rộng, khu vườn mới có:
	\begin{itemize}
		\item Chiều dài là: $70 + 2x\text{ (m)}$.
		\item Chiều rộng là: $30 + 2x\text{ (m)}$ (với $x > 0$).
	\end{itemize}
	Biểu thức biểu diễn diện tích khu vườn sau khi mở rộng là:
	$S = (70 + 2x)(30 + 2x) = 2100 + 140x + 60x + 4x^2 = 4x^2 + 200x + 2100\text{ (m}^2\text{)}$.
	\item Diện tích ban đầu của khu vườn là: $70\cdot 30 = 2100\text{ (m}^2\text{)}$.\\
	Theo đề bài, sau khi mở rộng thì diện tích khu vườn tăng thêm $1150\text{ m}^2$, nên ta có phương trình:
	$S - 2100 = 1150 \Leftrightarrow 4x^2 + 200x = 1150 \Leftrightarrow 4x^2 + 200x - 1150 = 0 \Leftrightarrow 2x^2 + 100x - 575 = 0$.\\
	Ta có $\Delta' = 50^2 - 2\cdot(-575) = 2500 + 1150 = 3650 > 0$.\\
	Vì $x > 0$ nên phương trình có nghiệm dương duy nhất là:
	$x = \dfrac{-50 + \sqrt{3650}}{2} \approx \dfrac{-50 + 60{,}415}{2} = \dfrac{10{,}415}{2} \approx 5{,}2\text{ (m)}$.\\
	Vậy giá trị của $x$ xấp xỉ bằng $5{,}2\text{ mét}$.
\end{enumerate}
}
\end{bt}

\begin{bt}%[TS10 4.2]%[ma:02fe1f3b]
% nguon: Đề tham khảo Tuyển sinh 10 TP.HCM 2025-2026 – Quận 1 – Đề 1 – Bài 4
Một chiếc bàn tròn có bán kính là $2\text{ m}$. Người ta muốn thiết kế một chiếc khăn trải bàn hình tròn có bán kính rộng hơn bán kính của mặt bàn là $x\text{ (m)}$ ($x > 0$).
\begin{enumerate}[a)]
	\item Viết biểu thức $S$ biểu diễn diện tích của tấm khăn trải bàn theo $x$.
	\item Biết rằng phần diện tích khăn rủ xuống ngoài mép bàn là $12{,}0576\text{ m}^2$. Tính giá trị của $x$ (lấy $\pi \approx 3{,}14$).
\end{enumerate}
\loigiai{
\begin{enumerate}[a)]
	\item Bán kính của tấm khăn trải bàn hình tròn là: $R = 2 + x\text{ (m)}$.\\
	Biểu thức biểu diễn diện tích của tấm khăn trải bàn là:
	$S = \pi R^2 = \pi (2 + x)^2 = \pi (x^2 + 4x + 4)\text{ (m}^2\text{)}$.
	\item Diện tích mặt bàn tròn là: $S_{\text{bàn}} = \pi \cdot 2^2 = 4\pi\text{ (m}^2\text{)}$.\\
	Phần diện tích khăn rủ xuống ngoài mép bàn là:
	$S_{\text{dư}} = S - S_{\text{bàn}} = \pi (x^2 + 4x + 4) - 4\pi = \pi (x^2 + 4x)\text{ (m}^2\text{)}$.\\
	Theo giả thiết, $S_{\text{dư}} = 12{,}0576\text{ m}^2$ và $\pi \approx 3{,}14$, ta có:
	$3{,}14(x^2 + 4x) = 12{,}0576 \Leftrightarrow x^2 + 4x = 3{,}84 \Leftrightarrow x^2 + 4x - 3{,}84 = 0$.\\
	Ta có $\Delta' = 2^2 - 1\cdot(-3{,}84) = 4 + 3{,}84 = 7{,}84 = (2{,}8)^2 > 0$.\\
	Vì $x > 0$ nên phương trình có nghiệm dương duy nhất:
	$x = -2 + 2{,}8 = 0{,}8\text{ (m)}$.\\
	Vậy giá trị của $x$ là $0{,}8\text{ mét}$.
\end{enumerate}
}
\end{bt}

\begin{bt}%[TS10 4.1]%[ma:7f252767]
% nguon: Đề tham khảo Tuyển sinh 10 TP.HCM 2025-2026 – Quận 1 – Đề 2 – Bài 4
Một mảnh đất hình chữ nhật có chiều dài là $3x + 10\text{ (m)}$ và chiều rộng là $2x + 5\text{ (m)}$. Ở giữa mảnh đất người ta đào một cái ao hình tam giác đều có cạnh là $x + 1\text{ (m)}$, phần đất còn lại dùng để trồng rau.
\begin{enumerate}[a)]
	\item Viết biểu thức $A$ biểu thị diện tích phần đất trồng rau.
	\item Nếu diện tích của cái ao là $12\sqrt{3}\text{ m}^2$ thì diện tích phần trồng rau là bao nhiêu mét vuông (làm tròn kết quả đến hàng đơn vị)?
\end{enumerate}
\loigiai{
\begin{enumerate}[a)]
	\item Diện tích mảnh đất hình chữ nhật là:
	$S_{\text{đất}} = (3x + 10)(2x + 5) = 6x^2 + 15x + 20x + 50 = 6x^2 + 35x + 50\text{ (m}^2\text{)}$.\\
	Diện tích cái ao hình tam giác đều cạnh $x + 1$ là:
	$S_{\text{ao}} = \dfrac{\sqrt{3}}{4}(x + 1)^2\text{ (m}^2\text{)}$.\\
	Biểu thức $A$ biểu thị diện tích phần đất trồng rau là:
	$A = S_{\text{đất}} - S_{\text{ao}} = 6x^2 + 35x + 50 - \dfrac{\sqrt{3}}{4}(x + 1)^2\text{ (m}^2\text{)}$.
	\item Theo đề bài, diện tích cái ao là $12\sqrt{3}\text{ m}^2$, nên ta có:
	$\dfrac{\sqrt{3}}{4}(x + 1)^2 = 12\sqrt{3} \Leftrightarrow (x + 1)^2 = 48 \Leftrightarrow x + 1 = \sqrt{48} = 4\sqrt{3} \Leftrightarrow x = 4\sqrt{3} - 1 \approx 5{,}928\text{ (m)}$.\\
	Diện tích mảnh đất hình chữ nhật lúc này là:
	$S_{\text{đất}} = 6(4\sqrt{3} - 1)^2 + 35(4\sqrt{3} - 1) + 50 = 6(49 - 8\sqrt{3}) + 140\sqrt{3} - 35 + 50 = 294 - 48\sqrt{3} + 140\sqrt{3} + 15 = 309 + 92\sqrt{3}\text{ (m}^2\text{)}$.\\
	Diện tích phần trồng rau là:
	$A = S_{\text{đất}} - S_{\text{ao}} = (309 + 92\sqrt{3}) - 12\sqrt{3} = 309 + 80\sqrt{3} \approx 309 + 138{,}564 = 447{,}564\text{ (m}^2\text{)}$.\\
	Làm tròn đến hàng đơn vị, diện tích phần trồng rau xấp xỉ bằng $448\text{ m}^2$.
\end{enumerate}
}
\end{bt}

\begin{bt}%[TS10 4.1]%[ma:47256e5e]
% nguon: Đề tham khảo Tuyển sinh 10 TP.HCM 2025-2026 – Quận 5 – Đề 1 – Bài 4
Nhà ông Hiền có một mảnh đất hình chữ nhật có chiều rộng là $5\text{ m}$ và chiều dài là $50\text{ m}$. Do việc mở rộng đường, nên nhà nước đã phải thu hồi một phần diện tích đất của nhà ông Hiền (phần hình tam giác vuông ở góc mảnh đất, với $0 < x < 5$).
\begin{enumerate}[a)]
    \item Viết và thu gọn biểu thức $A$ biểu thị theo $x$ diện tích phần đất bị thu hồi của nhà ông Hiền, biết hai cạnh góc vuông của phần tam giác bị thu hồi lần lượt là $5 - x$ (m) và $50 - 2x$ (m).
    \item Phần đất bị thu hồi của nhà ông Hiền được nhà nước đền bù số tiền $440$ triệu đồng. Tìm giá trị $x$, biết giá đền bù đất bị thu hồi là $10\text{ triệu đồng/m}^2$.
\end{enumerate}
\loigiai{
\begin{enumerate}[a)]
    \item Phần đất bị thu hồi là một hình tam giác vuông có độ dài hai cạnh góc vuông lần lượt là $5 - x\text{ (m)}$ và $50 - 2x\text{ (m)}$.\\
    Diện tích phần đất bị thu hồi được biểu diễn bởi:
    $A = \dfrac{1}{2}(5 - x)(50 - 2x) = (5 - x)(25 - x) = x^2 - 30x + 125\text{ (m}^2\text{)}$.
    \item Diện tích phần đất bị thu hồi tương ứng với số tiền đền bù là:
    $S = \dfrac{440}{10} = 44\text{ (m}^2\text{)}$.\\
    Theo câu a, ta có phương trình:
    $x^2 - 30x + 125 = 44 \Leftrightarrow x^2 - 30x + 81 = 0 \Leftrightarrow (x - 3)(x - 27) = 0 \Leftrightarrow \left[\begin{aligned} &x = 3\text{ (thỏa mãn } 0 < x < 5) \\ &x = 27\text{ (loại vì } x > 5). \end{aligned}\right.$\\
    Vậy giá trị cần tìm là $x = 3$.
\end{enumerate}
}
\end{bt}

\begin{bt}%[TS10 4.3]%[ma:db59061f]
% nguon: Đề tham khảo Tuyển sinh 10 TP.HCM 2025-2026 – Quận 5 – Đề 3 – Bài 4
Chợ Dân Sinh chuẩn bị được đưa vào hoạt động. Nếu mỗi gian hàng của chợ này cho thuê với mức giá 31 triệu đồng/năm thì sẽ có 60 gian hàng được thuê. Theo khảo sát và nghiên cứu, nếu chợ giảm giá thuê và ước tính cứ mỗi lần giảm giá thuê 1 triệu đồng/năm thì số lượng gian hàng được thuê sẽ tăng thêm 20 gian hàng. Gọi $x$ (triệu đồng) là giá cho thuê mới ($0 < x < 31$).
\begin{enumerate}[a)]
    \item Viết biểu thức $P(x)$ biểu diễn theo $x$ doanh thu của chợ Dân Sinh.
    \item Phí duy trì mỗi gian hàng là 27 triệu đồng/năm. Hỏi giá cho thuê mới của chợ là bao nhiêu để lợi nhuận thu được là 245 triệu đồng? Biết lợi nhuận bằng doanh thu trừ đi tổng chi phí duy trì gian hàng.
\end{enumerate}
\loigiai{
\begin{enumerate}[a)]
    \item Khi giá cho thuê mới là $x$ (triệu đồng/năm), mức giảm giá so với giá ban đầu là:
    $31 - x$ (triệu đồng/năm).\\
    Cứ giảm 1 triệu đồng thì số gian hàng thuê tăng thêm 20 gian hàng, nên số gian hàng được thuê là:
    $N = 60 + 20(31 - x) = 60 + 620 - 20x = 680 - 20x$ (gian hàng).\\
    Doanh thu của chợ Dân Sinh theo giá thuê $x$ được biểu diễn bởi:
    $P(x) = x(680 - 20x) = -20x^2 + 680x$ (triệu đồng).
    \item Tổng chi phí duy trì của tất cả các gian hàng được thuê là:
    $C = 27(680 - 20x)$ (triệu đồng).\\
    Lợi nhuận của chợ được tính theo công thức:
    $\text{Lợi nhuận} = P(x) - C = x(680 - 20x) - 27(680 - 20x) = (x - 27)(680 - 20x)$.\\
    Để lợi nhuận đạt 245 triệu đồng, ta có phương trình:
    $\begin{aligned}
    (x - 27)(680 - 20x) = 245 &\Leftrightarrow -20x^2 + 1220x - 18360 = 245 \\
    &\Leftrightarrow -20x^2 + 1220x - 18605 = 0 \\
    &\Leftrightarrow 4x^2 - 244x + 3721 = 0 \\
    &\Leftrightarrow (2x - 61)^2 = 0 \Leftrightarrow 2x = 61 \Leftrightarrow x = 30{,}5\text{ (thỏa mãn } 0 < x < 31).
    \end{aligned}$\\
    Vậy giá cho thuê mới của mỗi gian hàng là $30{,}5$ triệu đồng/năm.
\end{enumerate}
}
\end{bt}

\begin{bt}%[TS10 4.1]%[ma:03aeb78f]
% nguon: Đề tham khảo Tuyển sinh 10 TP.HCM 2025-2026 – Quận 6 – Đề 2 – Bài 4
Một mảnh vườn hình chữ nhật có chiều rộng bằng $a\text{ (m)}$, chiều dài hơn chiều rộng $6\text{ (m)}$. Người ta làm lối đi xung quanh rộng $x\text{ (m)}$ (với $0 < 2x < a$).
\begin{enumerate}[a)]
    \item Viết biểu thức biểu thị diện tích phần còn lại của mảnh vườn theo $a$ và $x$.
    \item Tính diện tích phần còn lại khi $a = 30\text{ m}$, $x = 1\text{ m}$.
\end{enumerate}
\loigiai{
\begin{enumerate}[a)]
    \item Chiều rộng ban đầu của mảnh vườn là $a\text{ (m)}$.\\
    Chiều dài ban đầu của mảnh vườn là:
    $a + 6\text{ (m)}$.\\
    Vì lối đi xung quanh vườn có chiều rộng là $x\text{ (m)}$ nên:
    \begin{itemize}
        \item Chiều rộng phần đất còn lại là: $a - 2x\text{ (m)}$.
        \item Chiều dài phần đất còn lại là: $(a + 6) - 2x = a - 2x + 6\text{ (m)}$.
    \end{itemize}
    Biểu thức biểu thị diện tích phần đất còn lại của mảnh vườn là:
    $S = (a - 2x)(a - 2x + 6)\text{ (m}^2\text{)}$.
    \item Khi $a = 30\text{ m}$ và $x = 1\text{ m}$:\\
    Chiều rộng phần còn lại là:
    $30 - 2 \cdot 1 = 28\text{ (m)}$.\\
    Chiều dài phần còn lại là:
    $30 + 6 - 2 \cdot 1 = 34\text{ (m)}$.\\
    Diện tích phần còn lại của mảnh vườn là:
    $S = 28 \cdot 34 = 952\text{ (m}^2\text{)}$.
\end{enumerate}
}
\end{bt}

\begin{bt}%[TS10 4.1]%[ma:529822e0]
% nguon: Đề tham khảo Tuyển sinh 10 TP.HCM 2025-2026 – Quận 7 – Đề 1 – Bài 4
Người ta làm một lối đi theo chiều dài và chiều rộng của một sân cỏ hình chữ nhật. Em hãy tính chiều rộng $x$ của lối đi, biết rằng lối đi có diện tích bằng $46\text{ m}^2$, sân cỏ ban đầu có chiều dài $15\text{ m}$ và chiều rộng $6\text{ m}$ (với $0 < x < 6$).
\begin{enumerate}[a)]
    \item Viết biểu thức tính diện tích của lối đi theo $x$.
    \item Tìm $x$.
\end{enumerate}
\loigiai{
\begin{enumerate}[a)]
    \item Diện tích toàn bộ sân cỏ ban đầu là:
    $S_0 = 15 \cdot 6 = 90\text{ (m}^2\text{)}$.\\
    Khi bớt lối đi rộng $x\text{ (m)}$ dọc theo một chiều dài và một chiều rộng của sân:
    \begin{itemize}
        \item Chiều dài phần sân cỏ còn lại là: $15 - x\text{ (m)}$.
        \item Chiều rộng phần sân cỏ còn lại là: $6 - x\text{ (m)}$.
    \end{itemize}
    Diện tích phần sân cỏ còn lại là:
    $S_{\text{còn}} = (15 - x)(6 - x) = x^2 - 21x + 90\text{ (m}^2\text{)}$.\\
    Diện tích của lối đi là hiệu diện tích sân ban đầu và phần còn lại:
    $S_{\text{lối đi}} = S_0 - S_{\text{còn}} = 90 - (x^2 - 21x + 90) = 21x - x^2\text{ (m}^2\text{)}$.
    \item Theo đề bài, diện tích lối đi bằng $46\text{ m}^2$, ta có phương trình:
    $21x - x^2 = 46 \Leftrightarrow x^2 - 21x + 46 = 0 \Leftrightarrow (x - 2)(x - 23) = 0 \Leftrightarrow \left[\begin{aligned} &x = 2\text{ (thỏa mãn } 0 < x < 6) \\ &x = 23\text{ (loại vì } x > 6). \end{aligned}\right.$\\
    Vậy chiều rộng của lối đi là $x = 2\text{ m}$.
\end{enumerate}
}
\end{bt}

\begin{bt}%[TS10 4.2]%[ma:478402a9]
% nguon: Đề tham khảo Tuyển sinh 10 TP.HCM 2025-2026 – Quận 8 – Đề 1 – Câu 4
Một cái sân hình chữ nhật có chiều dài là $1{,}57x + 10\text{ (m)}$ và chiều rộng là $2x + 2\text{ (m)}$ (với $x > 0$). Ở góc sân, người ta làm một bồn hoa hình tròn có bán kính $x\text{ (m)}$. (Cho $\pi = 3{,}14$).
\begin{enumerate}[a)]
    \item Viết biểu thức biểu thị diện tích đất còn lại của sân sau khi đã xây bồn hoa.
    \item Hãy tính bán kính $x$ của bồn hoa hình tròn, biết diện tích đất còn lại sau khi xây bồn hoa là $54{,}71\text{ m}^2$.
\end{enumerate}
\loigiai{
\begin{enumerate}[a)]
    \item Diện tích của toàn bộ cái sân hình chữ nhật là:
    $S_{\text{sân}} = (1{,}57x + 10)(2x + 2) = 3{,}14x^2 + 3{,}14x + 20x + 20 = 3{,}14x^2 + 23{,}14x + 20\text{ (m}^2\text{)}$.\\
    Diện tích bồn hoa hình tròn bán kính $x$ là:
    $S_{\text{hoa}} = \pi x^2 = 3{,}14x^2\text{ (m}^2\text{)}$.\\
    Biểu thức diện tích đất còn lại sau khi xây bồn hoa là:
    $S = S_{\text{sân}} - S_{\text{hoa}} = (3{,}14x^2 + 23{,}14x + 20) - 3{,}14x^2 = 23{,}14x + 20\text{ (m}^2\text{)}$.
    \item Theo đề bài, diện tích đất còn lại là $54{,}71\text{ m}^2$, ta có phương trình:
    $23{,}14x + 20 = 54{,}71 \Leftrightarrow 23{,}14x = 34{,}71 \Leftrightarrow x = \dfrac{34{,}71}{23{,}14} = 1{,}5\text{ (m)}$.\\
    Vậy bán kính của bồn hoa hình tròn là $x = 1{,}5\text{ m}$.
\end{enumerate}
}
\end{bt}

\begin{bt}%[TS10 4.1]%[ma:ce90aa05]
% nguon: Đề tham khảo Tuyển sinh 10 TP.HCM 2025-2026 – Quận 8 – Đề 2 – Bài 4
Trên một miếng đất hình chữ nhật có chiều dài là $(2x+10)\text{ m}$ và chiều rộng là $x\text{ m}$ ($x > 0$), người ta làm một lối đi hình bình hành có đáy nằm trên chiều dài với bề rộng là $2\text{ m}$. Phần còn lại là hai miếng đất hình thang vuông có diện tích bằng nhau dùng để trồng hoa hồng và hoa cúc.
\begin{enumerate}[a)]
    \item Viết biểu thức $A$ biểu diễn theo $x$ diện tích phần đất trồng hoa hồng.
    \item Người ta dự kiến lát sỏi lối đi với chi phí mỗi mét vuông là $120$ nghìn đồng. Hỏi tổng chi phí làm lối đi là bao nhiêu? Biết diện tích miếng đất trồng hoa hồng là $45\text{ m}^2$.
\end{enumerate}
\loigiai{
\begin{enumerate}[a)]
    \item Tổng diện tích miếng đất hình chữ nhật ban đầu là:
    $S_{\text{đất}} = x(2x + 10) = 2x^2 + 10x\text{ (m}^2\text{)}$.\\
    Lối đi hình bình hành có đáy bằng $2\text{ m}$ và chiều cao tương ứng chính là chiều rộng $x\text{ (m)}$ của mảnh đất hình chữ nhật.\\
    Diện tích lối đi là:
    $S_{\text{lối đi}} = 2 \cdot x = 2x\text{ (m}^2\text{)}$.\\
    Tổng diện tích đất còn lại để trồng hoa hồng và hoa cúc là:
    $S_{\text{còn}} = S_{\text{đất}} - S_{\text{lối đi}} = (2x^2 + 10x) - 2x = 2x^2 + 8x\text{ (m}^2\text{)}$.\\
    Vì hai phần đất trồng hoa hồng và hoa cúc có diện tích bằng nhau nên diện tích trồng hoa hồng là:
    $A = \dfrac{2x^2 + 8x}{2} = x^2 + 4x\text{ (m}^2\text{)}$.
    \item Theo đề bài, diện tích trồng hoa hồng là $45\text{ m}^2$, ta có phương trình:
    $x^2 + 4x = 45 \Leftrightarrow x^2 + 4x - 45 = 0 \Leftrightarrow (x - 5)(x + 9) = 0 \Leftrightarrow \left[\begin{aligned} &x = 5\text{ (thỏa mãn } x > 0) \\ &x = -9\text{ (loại)}. \end{aligned}\right.$\\
    Với $x = 5\text{ m}$, diện tích lối đi là:
    $S_{\text{lối đi}} = 2 \cdot 5 = 10\text{ (m}^2\text{)}$.\\
    Chi phí để lát sỏi làm lối đi là:
    $10 \cdot 120\,000 = 1\,200\,000\text{ (đồng)}$.\\
    Vậy chi phí để làm lối đi là $1\,200\,000$ đồng ($1{,}2$ triệu đồng).
\end{enumerate}
}
\end{bt}

\begin{bt}%[TS10 4.1]%[ma:a11428f8]
% nguon: Đề tham khảo Tuyển sinh 10 năm học 2025-2026 Sở GD&ĐT TP.HCM – Quận 10 - Đề 2 (Bài 3)
Một chiếc hộp không có nắp được làm từ một mảnh bìa các-tông. Chiếc hộp có đáy là một hình vuông cạnh $x$ (cm) và chiều cao là $h = 10\text{ cm}$.
\begin{enumerate}[a)]
    \item Viết biểu thức biểu diễn diện tích xung quanh $S_{\text{xq}}$ của chiếc hộp theo $x$.
    \item Biết rằng tổng diện tích xung quanh và diện tích đáy của chiếc hộp là $225\text{ cm}^2$. Tính độ dài cạnh đáy của chiếc hộp.
\end{enumerate}
\loigiai{
\begin{enumerate}[a)]
    \item Chiếc hộp có đáy là hình vuông cạnh $x$ (cm), bốn mặt bên là bốn hình chữ nhật có kích thước $x\text{ cm} \times 10\text{ cm}$.
    Do đó diện tích xung quanh của chiếc hộp là:
    $S_{\text{xq}} = 4 \cdot (x \cdot 10) = 40x\text{ (cm}^2\text{)}$.
    \item Diện tích đáy của chiếc hộp hình vuông là:
    $S_{\text{đáy}} = x^2\text{ (cm}^2\text{)}$ (với $x > 0$).
    Theo đề bài, tổng diện tích xung quanh và diện tích đáy bằng $225\text{ cm}^2$, ta có phương trình:
    $x^2 + 40x = 225 \iff x^2 + 40x - 225 = 0$.
    Ta có $\Delta' = 20^2 - 1 \cdot (-225) = 400 + 225 = 625 = 25^2 > 0$.
    Phương trình có hai nghiệm:
    $x_1 = -20 + 25 = 5$ (nhận vì $x > 0$),
    $x_2 = -20 - 25 = -45$ (loại).
    Vậy độ dài cạnh đáy của chiếc hộp là $5\text{ cm}$.
\end{enumerate}
}
\end{bt}

\begin{bt}%[TS10 4.1]%[ma:d14b5a28]
% nguon: Đề tham khảo Tuyển sinh 10 năm học 2025-2026 Sở GD&ĐT TP.HCM – Quận 11 - Đề 2 (Bài 3)
Một sân chơi hình chữ nhật có chiều dài là $x$ mét và chiều dài hơn chiều rộng là $3\text{ m}$.
\begin{enumerate}[a)]
    \item Viết biểu thức tính diện tích $S$ của sân chơi hình chữ nhật theo $x$.
    \item Sau đó, ban quản lý quyết định tăng chiều dài thêm $1\text{ m}$ và tăng chiều rộng thêm $3\text{ m}$ thì diện tích của sân chơi hình chữ nhật tăng gấp đôi. Tính chiều dài ban đầu của sân chơi đó.
\end{enumerate}
\loigiai{
\begin{enumerate}[a)]
    \item Chiều rộng của sân chơi hình chữ nhật ban đầu là $x - 3$ (m), điều kiện $x > 3$.
    Diện tích ban đầu của sân chơi là:
    $S = x(x - 3) = x^2 - 3x\text{ (m}^2\text{)}$.
    \item Chiều dài mới sau khi tăng thêm $1\text{ m}$ là: $x + 1$ (m).
    Chiều rộng mới sau khi tăng thêm $3\text{ m}$ là: $(x - 3) + 3 = x$ (m).
    Diện tích mới của sân chơi là: $(x + 1)x = x^2 + x\text{ (m}^2\text{)}$.
    Vì diện tích mới tăng gấp đôi diện tích ban đầu nên ta có phương trình:
    $x^2 + x = 2(x^2 - 3x) \iff x^2 + x = 2x^2 - 6x \iff x^2 - 7x = 0 \iff x(x - 7) = 0$.
    Vì $x > 3$ nên $x = 7$ (thỏa mãn điều kiện).
    Vậy chiều dài ban đầu của sân chơi là $7\text{ m}$.
\end{enumerate}
}
\end{bt}

\begin{bt}%[TS10 4.1]%[ma:7a6973e2]
% nguon: Đề tham khảo Tuyển sinh 10 năm học 2025-2026 Sở GD&ĐT TP.HCM – Quận 12 - Đề 1 (Bài 4)
Một mảnh vườn hình chữ nhật có chiều rộng bằng $a$ (m) và chiều dài hơn chiều rộng $6\text{ m}$. Người ta bớt ra hai dải đất chạy dọc theo mép mảnh vườn có bề rộng là $x$ (m) để làm lối đi (với $0 < x < a$).
\begin{enumerate}[a)]
    \item Viết biểu thức biểu thị diện tích phần đất còn lại của mảnh vườn để trồng trọt theo $a$ và $x$.
    \item Tính diện tích phần đất còn lại khi $a = 30\text{ m}$ và $x = 1\text{ m}$.
\end{enumerate}
\loigiai{
\begin{enumerate}[a)]
    \item Chiều rộng ban đầu của mảnh vườn là $a$ (m), chiều dài ban đầu là $a + 6$ (m).
    Sau khi bớt lối đi rộng $x$ (m), kích thước phần đất còn lại là:
    Chiều rộng còn lại: $a - x$ (m);
    Chiều dài còn lại: $(a + 6) - x$ (m).
    Diện tích phần đất còn lại của mảnh vườn là:
    $S = (a - x)(a + 6 - x)\text{ (m}^2\text{)}$.
    \item Khi $a = 30\text{ m}$ và $x = 1\text{ m}$, ta có:
    Chiều rộng còn lại: $30 - 1 = 29\text{ m}$.
    Chiều dài còn lại: $(30 + 6) - 1 = 35\text{ m}$.
    Diện tích phần đất còn lại là:
    $S = 29 \cdot 35 = 1015\text{ m}^2$.
\end{enumerate}
}
\end{bt}

\begin{bt}%[TS10 4.1]%[ma:941757cb]
% nguon: Đề tham khảo Tuyển sinh 10 năm học 2025-2026 Sở GD&ĐT TP.HCM – Quận 12 - Đề 2 (Bài 3)
Bác Lâm sở hữu một khu đất hình thang có đáy nhỏ dài $15\text{ m}$, đáy lớn dài $25\text{ m}$ và chiều cao bằng $18\text{ m}$. Bác dự định làm một lối đi có dạng hình bình hành xuyên qua khu đất với chiều rộng đáy bằng $x$ (m) ($0 < x < 15$) và chiều cao tương ứng là $18\text{ m}$. Toàn bộ phần diện tích còn lại của khu đất bác Lâm dùng để trồng hoa hồng.
\begin{enumerate}[a)]
    \item Viết biểu thức biểu diễn diện tích đất trồng hoa hồng theo $x$.
    \item Hãy tính độ rộng $x$ của lối đi biết diện tích đất trồng hoa hồng là $324\text{ m}^2$.
\end{enumerate}
\loigiai{
\begin{enumerate}[a)]
    \item Diện tích của toàn bộ khu đất hình thang là:
    $S_{\text{khu đất}} = \dfrac{(15 + 25) \cdot 18}{2} = 40 \cdot 9 = 360\text{ (m}^2\text{)}$.
    Lối đi hình bình hành có cạnh đáy là $x$ (m) và chiều cao tương ứng $18\text{ m}$ nên có diện tích là:
    $S_{\text{lối đi}} = 18x\text{ (m}^2\text{)}$.
    Diện tích phần đất còn lại để trồng hoa hồng là:
    $S_{\text{hoa}} = S_{\text{khu đất}} - S_{\text{lối đi}} = 360 - 18x\text{ (m}^2\text{)}$.
    \item Khi diện tích trồng hoa hồng bằng $324\text{ m}^2$, ta có phương trình:
    $360 - 18x = 324 \iff 18x = 36 \iff x = 2$ (m) (thỏa mãn điều kiện $0 < x < 15$).
    Vậy chiều rộng của lối đi là $2\text{ m}$.
\end{enumerate}
}
\end{bt}

\begin{bt}%[TS10 4.1]%[ma:430d5618]
% nguon: Đề tham khảo Tuyển sinh 10 năm học 2025-2026 Sở GD&ĐT TP.HCM – TP. Thủ Đức - Đề 1 (Bài 4)
Một mảnh đất vườn có dạng hình vuông với độ dài cạnh là $x$ mét ($x > 0$). Chủ vườn mở rộng mảnh đất thành một hình chữ nhật có chiều rộng giữ nguyên là $x$ (m) và chiều dài tăng thêm $3\text{ m}$ (tức chiều dài mới là $x + 3\text{ m}$).
\begin{enumerate}[a)]
    \item Viết biểu thức rút gọn $S$ biểu diễn phần diện tích tăng thêm của mảnh vườn theo $x$.
    \item Tìm cạnh $x$ của mảnh vườn ban đầu, biết rằng diện tích tăng thêm của mảnh vườn đúng bằng $21\text{ m}^2$.
\end{enumerate}
\loigiai{
\begin{enumerate}[a)]
    \item Diện tích mảnh đất hình vuông ban đầu là: $S_1 = x^2\text{ (m}^2\text{)}$.
    Diện tích mảnh đất hình chữ nhật sau khi mở rộng là:
    $S_2 = x(x + 3) = x^2 + 3x\text{ (m}^2\text{)}$.
    Phần diện tích tăng thêm của mảnh vườn là:
    $S = S_2 - S_1 = (x^2 + 3x) - x^2 = 3x\text{ (m}^2\text{)}$.
    \item Theo đề bài, phần diện tích tăng thêm là $21\text{ m}^2$, ta có phương trình:
    $3x = 21 \iff x = 7$ (m) (thỏa mãn $x > 0$).
    Vậy độ dài cạnh của mảnh vườn hình vuông ban đầu là $7\text{ m}$.
\end{enumerate}
}
\end{bt}

\begin{bt}%[TS10 4.1]%[ma:e46e249c]
% nguon: Đề tham khảo Tuyển sinh 10 năm học 2025-2026 Sở GD&ĐT TP.HCM – TP. Thủ Đức - Đề 2 (Bài 4)
Người ta làm một lối đi bao quanh dọc theo một chiều dài và một chiều rộng của một hồ nước hình chữ nhật có kích thước dài $6\text{ m}$ và rộng $5\text{ m}$ (lối đi có bề rộng đồng đều là $x$ mét).
\begin{enumerate}[a)]
    \item Viết biểu thức biểu diễn diện tích $S$ của phần lối đi theo $x$.
    \item Tìm độ rộng $x$ của lối đi biết diện tích của lối đi đó bằng $50\text{ m}^2$ (kết quả làm tròn đến chữ số thập phân thứ hai).
\end{enumerate}
\loigiai{
\begin{enumerate}[a)]
    \item Diện tích của hồ nước hình chữ nhật ban đầu là:
    $S_{\text{hồ}} = 6 \cdot 5 = 30\text{ (m}^2\text{)}$.
    Khu đất bao gồm cả hồ nước và phần lối đi mở rộng có kích thước là $6 + x$ (m) và $5 + x$ (m).
    Diện tích của toàn bộ khu đất mới là:
    $S_{\text{tổng}} = (6 + x)(5 + x) = x^2 + 11x + 30\text{ (m}^2\text{)}$.
    Biểu thức diện tích phần lối đi là:
    $S = S_{\text{tổng}} - S_{\text{hồ}} = (x^2 + 11x + 30) - 30 = x^2 + 11x\text{ (m}^2\text{)}$.
    \item Khi diện tích lối đi bằng $50\text{ m}^2$, ta có phương trình:
    $x^2 + 11x = 50 \iff x^2 + 11x - 50 = 0$.
    Ta có $\Delta = 11^2 - 4 \cdot 1 \cdot (-50) = 121 + 200 = 321 > 0$.
    Vì độ rộng $x > 0$ nên nghiệm thỏa mãn là:
    $x = \dfrac{-11 + \sqrt{321}}{2} \approx 3{,}46\text{ m}$.
    Vậy độ rộng của lối đi xấp xỉ $3{,}46\text{ m}$.
\end{enumerate}
}
\end{bt}
